\Chapter{Grafos dirigidos y flujo en redes}

Estas notas presentan, con el rigor propio de la matemática pura, la teoría que subyace al problema clásico de la asignación de tráfico en redes viales. El punto de partida es la observación, aparentemente simple, de que en una red vial los conductores eligen sus rutas de forma egoísta, buscando minimizar su propio tiempo de viaje. Esta elección descentralizada produce una distribución de flujos que, en general, no minimiza el costo total del sistema. La brecha entre ambos objetivos ---el individual y el colectivo--- es lo que Koutsoupias y Papadimitriou denominaron el \emph{Precio de la Anarquía}.

\section{Grafos dirigidos}

\subsection{Definiciones básicas}

\datebox{2026}

\begin{definition}[Grafo dirigido]
Un \emph{grafo dirigido} (o dígrafo) es un par ordenado $G = (V, A)$ donde:
\begin{itemize}
    \item $V$ es un conjunto finito no vacío, cuyos elementos se denominan \emph{nodos} o \emph{vértices};
    \item $A \subseteq V \times V \setminus \{(v,v) : v \in V\}$ es un conjunto finito de \emph{aristas dirigidas}.
\end{itemize}
Una arista $a = (u, v) \in A$ tiene \emph{cola} $\mathrm{tail}(a) = u$ y \emph{cabeza} $\mathrm{head}(a) = v$. Se dice que $a$ \emph{sale} de $u$ y \emph{entra} en $v$.
\end{definition}

\begin{notation}
Denotamos $|V| = n$ y $|A| = m$. Para un nodo $v \in V$, el \emph{conjunto de aristas salientes} es $\delta^+(v) \defeq \{a \in A : \mathrm{tail}(a) = v\}$ y el \emph{conjunto de aristas entrantes} es $\delta^-(v) \defeq \{a \in A : \mathrm{head}(a) = v\}$. El \emph{grado de salida} es $d^+(v) = |\delta^+(v)|$ y el \emph{grado de entrada} es $d^-(v) = |\delta^-(v)|$.
\end{notation}

\begin{definition}[Trayectoria]
Una \emph{trayectoria} (o camino dirigido) de $s$ a $d$ en $G$ es una secuencia alternada de nodos y aristas:
\[
p = (v_0, a_1, v_1, a_2, v_2, \ldots, a_k, v_k)
\]
tal que $v_0 = s$, $v_k = d$, y $a_i = (v_{i-1}, v_i) \in A$ para todo $i = 1, \ldots, k$. La trayectoria es \emph{simple} si no repite nodos: $v_i \neq v_j$ para $i \neq j$. Sea $\calP_{sd}$ el conjunto de todas las trayectorias simples de $s$ a $d$.
\end{definition}

\begin{remark}
En el contexto de redes de transporte, cuando se habla de ``rutas'' se hace referencia a trayectorias simples. La restricción de simplicidad es necesaria para que el problema de optimización esté bien planteado, ya que una trayectoria con ciclos tendría costo no acotado superiormente.
\end{remark}

\subsection{Flujo en redes}

\begin{definition}[Conjunto de pares origen-destino]
Un \emph{par origen-destino} (par O-D) es un par ordenado $w = (s_w, d_w) \in V \times V$ con $s_w \neq d_w$. Sea $W$ el conjunto finito de pares O-D. A cada par $w \in W$ se asocia una \emph{demanda} $d_w > 0$, que representa el flujo total que debe ser enviado de $s_w$ a $d_w$. Sea $\calP_w \defeq \calP_{s_w d_w}$ el conjunto de trayectorias simples para el par $w$, y $\calP \defeq \bigcup_{w \in W} \calP_w$.
\end{definition}

\begin{definition}[Vector de flujo]
Un \emph{vector de flujo por rutas} es un vector $f \in \R^{|\calP|}_+$ donde $f_p \geq 0$ denota el flujo asignado a la ruta $p \in \calP$. El flujo \emph{por aristas} inducido por $f$ es:
\[
x_a = x_a(f) \defeq \sum_{p \in \calP} \delta_{ap} f_p \quad \forall a \in A,
\]
donde $\delta_{ap} = 1$ si $a \in p$ y $\delta_{ap} = 0$ en caso contrario. En notación matricial, $x = \Delta f$, donde $\Delta \in \{0,1\}^{m \times |\calP|}$ es la \emph{matriz de incidencia arista-ruta}.
\end{definition}

\begin{definition}[Conjunto factible]
El conjunto de flujos factibles es:
\[
\FeasSet \defeq \left\{ f \in \R^{|\calP|}_+ \;\middle|\; \sum_{p \in \calP_w} f_p = d_w \quad \forall w \in W \right\}.
\]
\end{definition}

\begin{proposition}\label{prop:omega_compact_convex}
El conjunto $\FeasSet$ es compacto y convexo.
\end{proposition}

\begin{proof}
La convexidad es inmediata: si $f^1, f^2 \in \FeasSet$ y $\lambda \in [0,1]$, entonces $\lambda f^1 + (1-\lambda)f^2 \geq 0$ y
\[
\sum_{p \in \calP_w} [\lambda f^1_p + (1-\lambda) f^2_p] = \lambda d_w + (1-\lambda) d_w = d_w \quad \forall w \in W.
\]
Para la compacidad, observamos que $\FeasSet$ es cerrado (definido por igualdades lineales y desigualdades no estrictas) y acotado: para cada ruta $p \in \calP_w$, se tiene $0 \leq f_p \leq d_w$. Por el teorema de Heine-Borel, $\FeasSet$ es compacto.
\end{proof}

\subsection{La ley de conservación de flujo}

\begin{definition}[Conservación de flujo]
Dado un vector de flujo por aristas $x \in \R^m_+$, decimos que $x$ satisface la \emph{ley de conservación de flujo} para el par $w = (s,d)$ con demanda $d_w$ si:
\[
\sum_{a \in \delta^+(v)} x_a^w - \sum_{a \in \delta^-(v)} x_a^w = b_v^w \quad \forall v \in V,
\]
donde $b_v^w = d_w$ si $v = s$, $b_v^w = -d_w$ si $v = d$, y $b_v^w = 0$ en caso contrario.
\end{definition}

\begin{theorem}[Equivalencia de formulaciones]\label{thm:equivalencia}
Sea $x \in \R^m_+$. Las siguientes afirmaciones son equivalentes:
\begin{enumerate}
    \item Existe $f \in \FeasSet$ tal que $x = \Delta f$.
    \item $x$ satisface la ley de conservación de flujo para todo par $w \in W$.
\end{enumerate}
\end{theorem}

\begin{proof}
$(1) \Rightarrow (2)$: Si $x = \Delta f$ con $f \in \FeasSet$, entonces para cada nodo $v$:
\[
\sum_{a \in \delta^+(v)} x_a - \sum_{a \in \delta^-(v)} x_a = \sum_{a \in \delta^+(v)} \sum_p \delta_{ap} f_p - \sum_{a \in \delta^-(v)} \sum_p \delta_{ap} f_p.
\]
Para cada ruta $p \in \calP_w$ de $s$ a $d$, el número de aristas de $p$ que salen de $v$ menos el número que entran en $v$ es $\mathbbm{1}[v=s] - \mathbbm{1}[v=d]$. Por tanto la suma anterior vale $\sum_{p \in \calP_w} f_p (\mathbbm{1}[v=s_w] - \mathbbm{1}[v=d_w]) = b_v^w$.

$(2) \Rightarrow (1)$: Es el \emph{teorema de descomposición de flujo}: todo flujo que satisface conservación puede escribirse como combinación convexa de flujos de ruta. La demostración es constructiva: se identifica una trayectoria $p$ de $s$ a $d$ con flujo positivo en cada arista (posible mientras quede flujo no asignado), se le asigna el mínimo flujo de sus aristas, y se itera.
\end{proof}

\begin{remark}
El Teorema~\upref{thm:equivalencia} justifica trabajar indistintamente con variables de ruta o variables de arista. En el análisis teórico es conveniente usar variables de ruta; en los algoritmos, variables de arista.
\end{remark}

\section{Algoritmos de rutas más cortas}

\subsection{El algoritmo de Dijkstra}

\begin{definition}[Problema de la ruta más corta]
Dado $G = (V, A)$ con pesos $w : A \to \R_+$ y un nodo fuente $s \in V$, el \emph{problema de la ruta más corta} consiste en encontrar, para cada $v \in V$, una trayectoria de $s$ a $v$ que minimice la suma de pesos de sus aristas.
\end{definition}

\begin{algorithm}
\caption{Dijkstra (1959)}
\label{alg:dijkstra}
\begin{algorithmic}[1]
\Require Grafo $G=(V,A)$, pesos $w_a \geq 0$, nodo fuente $s$
\Ensure Distancias mínimas $d[v]$ para todo $v \in V$
\State $d[s] \leftarrow 0$; $d[v] \leftarrow \infty$ para todo $v \neq s$
\State $S \leftarrow \emptyset$; $Q \leftarrow V$ \Comment{Cola de prioridad}
\While{$Q \neq \emptyset$}
    \State $u \leftarrow \arg\min_{v \in Q} d[v]$
    \State $Q \leftarrow Q \setminus \{u\}$; $S \leftarrow S \cup \{u\}$
    \For{cada arista $(u,v) \in A$}
        \If{$d[u] + w_{(u,v)} < d[v]$}
            \State $d[v] \leftarrow d[u] + w_{(u,v)}$
        \EndIf
    \EndFor
\EndWhile
\end{algorithmic}
\end{algorithm}

\begin{theorem}[Correctitud de Dijkstra]
Al terminar el Algoritmo~\upref{alg:dijkstra}, $d[v]$ es la longitud de la ruta más corta de $s$ a $v$ para todo $v \in V$ alcanzable desde $s$.
\end{theorem}

\begin{proof}
Por inducción sobre $|S|$. Invariante: en cada iteración, $d[u] = \delta(s,u)$ (distancia mínima real) para todo $u \in S$.

\emph{Base}: $|S| = 0$. Vacuamente cierto.

\emph{Paso inductivo}: Supongamos que el invariante se cumple hasta extraer $k$ nodos. Sea $u$ el $(k+1)$-ésimo nodo extraído, con $d[u]$ en el momento de extracción. Supongamos por contradicción que existe una ruta $q$ de $s$ a $u$ con $\ell(q) < d[u]$. Sea $y$ el primer nodo de $q$ fuera de $S$, y $x$ su predecesor en $S$. Entonces:
\[
\ell(q) \geq d[x] + w_{(x,y)} = d[y] \geq d[u],
\]
donde la primera igualdad usa la hipótesis inductiva para $x$ y la definición de $d[y]$ en la actualización, y la segunda igualdad usa que $u$ fue elegido por tener $d[u] \leq d[y]$. Contradicción.
\end{proof}

\begin{proposition}[Complejidad de Dijkstra]
Con una cola de prioridad basada en montículo de Fibonacci, el Algoritmo~\upref{alg:dijkstra} tiene complejidad $O(n \log n + m)$.
\end{proposition}

\subsection{El algoritmo de Yen --- $k$ rutas más cortas}

\begin{definition}[Problema de las $k$ rutas más cortas]
Dado $G = (V,A)$, pesos $w_a \geq 0$, y un par $(s,d) \in V \times V$, el problema consiste en encontrar las $k$ trayectorias simples de $s$ a $d$ con menor costo, ordenadas de menor a mayor.
\end{definition}

\begin{algorithm}
\caption{Algoritmo de Yen (1971)}
\label{alg:yen}
\begin{algorithmic}[1]
\Require Grafo $G$, pesos $w$, nodos $s, d$, entero $k \geq 1$
\Ensure Las $k$ rutas simples más cortas de $s$ a $d$
\State $A_1 \leftarrow$ ruta más corta de $s$ a $d$ (Dijkstra)
\State $B \leftarrow \emptyset$ \Comment{Candidatos}
\For{$i = 2, \ldots, k$}
    \For{$j = 0, \ldots, |A_{i-1}| - 2$}
        \State $v_j \leftarrow j$-ésimo nodo de $A_{i-1}$
        \State $G' \leftarrow G$ con las aristas $(v_j, v_{j+1}^{(\ell)})$ eliminadas para cada ruta $A_\ell$ que comparte el prefijo $A_{i-1}[s \to v_j]$
        \State $p_{\mathrm{desv}} \leftarrow$ ruta más corta de $v_j$ a $d$ en $G'$
        \If{$p_{\mathrm{desv}} \neq \emptyset$}
            \State $B \leftarrow B \cup \{A_{i-1}[s \to v_j] \oplus p_{\mathrm{desv}}\}$
        \EndIf
    \EndFor
    \If{$B = \emptyset$}
        \State \textbf{break}
    \EndIf
    \State $A_i \leftarrow \arg\min_{p \in B} \ell(p)$; $B \leftarrow B \setminus \{A_i\}$
\EndFor
\end{algorithmic}
\end{algorithm}

\begin{theorem}[Correctitud de Yen]
El Algoritmo~\upref{alg:yen} genera correctamente las $k$ trayectorias simples más cortas de $s$ a $d$, sin repetición.
\end{theorem}

\begin{proof}[Esquema de la demostración]
La idea clave es que toda trayectoria simple $p$ de $s$ a $d$ distinta de $A_1, \ldots, A_{i-1}$ puede escribirse como $A_\ell[s \to v_j] \oplus q$ para alguna ruta anterior $A_\ell$ y algún nodo de desvío $v_j$, donde $q$ es una trayectoria de $v_j$ a $d$ que no usa los arcos de desvío de las rutas anteriores con el mismo prefijo. Por tanto, el candidato de menor costo en $B$ en cada iteración es necesariamente la $i$-ésima ruta más corta. La demostración por inducción sobre $i$ es estándar; véase \cite{Yen1971}.
\end{proof}

\begin{proposition}[Complejidad de Yen]
El Algoritmo~\upref{alg:yen} tiene complejidad $O(k n (m + n \log n))$.
\end{proposition}
